% !PRINTED 54
\weberpage{54}

\gesperrt{The number of genera is thus as great as the number of two-sided elements}\footnote{\gesperrt{Gauss} first introduced these concepts into the theory of quadratic forms, and used the expression ``Genera''. Disqu. arithm. art. 228 f.}.

\begin{center}
\S.\,14.
\end{center}

\begin{center}
\gesperrt{Indices for an odd prime power}\\
\gesperrt{as modulus.}
\end{center}

The most important example of a commutative group is offered by the natural numbers, when they are combined with one another by ordinary multiplication.

In order to derive from them a finite group, we take an arbitrary positive integer $m$ as modulus and decompose $m$ into its prime factors

\weq{1}{m = 2^\lambda q_1^{\varkappa_1} q_2^{\varkappa_2} \ldots,}

in which $q_1, q_2 \ldots$ denote distinct odd primes, $\lambda, \varkappa_1, \varkappa_2 \ldots$ positive exponents, $\lambda$ possibly also zero, namely when $m$ is odd.

Now all numbers, positive as well as negative, that are congruent to one another with respect to the modulus $m$ are united into a \gesperrt{number class}, and each of these number classes is represented by a representative, say by the residue of the division, hence by one of the numbers $0, 1, 2 \ldots m-1$.

If $a$ and $a'$ are two numbers of such a class, and $b$ and $b'$ two numbers of a second class, then the products $ab$ and $a'b'$ also belong to the same class. For if $a \equiv a'$, $b \equiv b'$, then also $ab \equiv a'b'\ (\mathrm{mod}\ m)$. By multiplication, therefore, not only the numbers but also the classes are compounded.

Nevertheless these number classes in their totality still form no group; for from a congruence

\[
ab \equiv ac\ (\mathrm{mod}\ m)
\]

it follows necessarily that $b \equiv c$ only when $a$ is relatively prime to $m$. Thus the requirement \S.\,1, 3. is here in general not fulfilled.

The greatest common divisor which a number $a$ has with $m$ is, for the entire class represented by $a$, the
%CONT>>
